\section{Keep the old and new working values separate}
The compression state contains eight words. At the beginning of a block, copy them
into working values called \(a,b,c,d,e,f,g,h\). During each round, compute two
temporary values from the \emph{old} working values:
\[
  \begin{aligned}
    T_1 & =h+\Sigma_1(e)+\operatorname{Ch}(e,f,g)+K_t+W_t, \\
    T_2 & =\Sigma_0(a)+\operatorname{Maj}(a,b,c).
  \end{aligned}
\]
All additions in the compression equations are modulo \(2^{32}\). The constants
\(K_t\) are fixed by the algorithm; \(W_t\) comes from this block's schedule.
The capital-sigma functions use only rotations:
\[
  \begin{aligned}
    \Sigma_0(x) & =\operatorname{ROTR}_2(x)\oplus\operatorname{ROTR}_{13}(x)
    \oplus\operatorname{ROTR}_{22}(x),                                       \\
    \Sigma_1(x) & =\operatorname{ROTR}_6(x)\oplus\operatorname{ROTR}_{11}(x)
    \oplus\operatorname{ROTR}_{25}(x).
  \end{aligned}
\]
Do not substitute the small-sigma functions from schedule expansion. Similar names
identify different operations.

The choice function operates independently at each bit position:
\(\operatorname{Ch}(e,f,g)=(e\mathbin{\&}f)\oplus((\mathord{\sim}e)\mathbin{\&}g)\).
A one bit in \(e\) selects the corresponding bit from \(f\); a zero bit selects it
from \(g\). An equivalent expression is \(g\oplus(e\mathbin{\&}(f\oplus g))\).
The majority function is
\(\operatorname{Maj}(a,b,c)=(a\mathbin{\&}b)\oplus(a\mathbin{\&}c)\oplus(b\mathbin{\&}c)\).
At a single bit position its result is one when at least two inputs are one.
The tilde denotes a bitwise complement within a word, and \(\&\) denotes bitwise AND.

After both temporaries are ready, the new working values are
\[
  (a',b',c',d',e',f',g',h')=
  (T_1+T_2,a,b,c,d+T_1,e,f,g).
\]
The primes distinguish the new round's values. With one register per working word,
move from the end of the list toward the beginning. Copy old \(g\) to \(h\) before
overwriting \(g\), and old \(f\) to \(g\) before overwriting \(f\). Reversing that
order can make several registers contain the same old value. The error is then carried
through every later round.

After all 64 rounds, add each working word to its corresponding incoming state word
and store the sum. This is the feed-forward step. The working values by themselves
are not the next compression state. Preserve the incoming words in memory until this
addition; loading them only after overwriting the state would lose the values needed.
